\documentclass[12pt]{article}
\usepackage[utf8]{inputenc}
\usepackage[T1]{fontenc}
\usepackage{hyperref}
\title{Uruava Language Resources for K-12 Curriculum Development}
\author{Compiled by Claude Sonnet 5 (claude-sonnet-5) for Andr\'as Kornai}
\date{}
\begin{document}
\maketitle

\section{Introduction}
Uruava (ISO 639-3 urv) is an extinct Austronesian language formerly spoken around Arawa on Bougainville, Papua New Guinea. It belongs to the Mono--Uruavan branch of Northwest Solomonic. One source reports that the last native speaker died around 1970. Speakers shifted to Nasioi. A revitalization curriculum would have to be built from historical documentation only.

\section{Primary Educational Resources}
I found no primers, textbooks or school materials for Uruava.

\section{Linguistic and Reference Materials}
\begin{itemize}
\item Uruava dictionary (Surrey Morphology Group, Northwest Solomonic pages): \url{https://www.smg.surrey.ac.uk/languages/northwest-solomonic/uruava/uruava-dictionary/v2}
\item Encyclopedic overview: \url{https://en.wikipedia.org/wiki/Uruava_language}
\item Capell's 1960s notes on the language, held in the PARADISEC archive (series NSBURU): \url{https://www.paradisec.org.au/fieldnotes/NSBURU.htm}
\item Verbix page on the language (dormant): \url{https://docs.verbix.com/EndangeredLanguages/Uruava}
\end{itemize}

\section{Religious and Community Resources}
I found no scripture translations or community organisations for Uruava in the search results.

\section{Digital and Audio Resources}
The PARADISEC archive listing above is the only archival lead, and I did not confirm whether it includes audio.

\section{Curriculum Development Context}
The language has no known speakers, and the record is limited to old word lists, notes and a dictionary. A curriculum could introduce the documented vocabulary as heritage content alongside Nasioi and Tok Pisin. Any revitalization work would need the Arawa community's involvement.
\end{document}